% =============================================================================
% Project : THE_PHILOSOPHY_OF_VIABILITY
% Book    : Book I — Ontology of Viability
% Book ID : VS.B01
%
% File    : 3_BACKMATTER/05_acronyms.tex
%
% Component:
%   List of Acronyms and Structural Codes
%
% Purpose:
%   Provides the official list of abbreviations and structural codes
%   used in Book I of The Philosophy of Viability.
%
% Editorial Policy:
%   1. Only abbreviations and codes established in the project structure
%      are included.
%   2. This file does not introduce new philosophical concepts.
%   3. Structural and logical codes are distinguished from ordinary
%      linguistic abbreviations.
%   4. Command definitions belonging to macros.tex are not duplicated here.
%   5. The list may be extended when new abbreviations are formally adopted.
%
% Dependencies:
%   Standard LaTeX commands only.
%
% =============================================================================


% -----------------------------------------------------------------------------
% LIST OF ACRONYMS AND STRUCTURAL CODES
% -----------------------------------------------------------------------------

\chapter*{List of Acronyms and Structural Codes}
\addcontentsline{toc}{chapter}{List of Acronyms and Structural Codes}

\markboth
    {List of Acronyms and Structural Codes}
    {List of Acronyms and Structural Codes}


This list records the abbreviations and structural codes used to identify
the project, its books, and its formal logical registers.

The entries are grouped according to their function. A structural or
registry code is not, by itself, a philosophical proposition, an
ontological commitment, or an indication of logical priority.


% =============================================================================
% 01 — PROJECT AND BOOK IDENTIFIERS
% =============================================================================

\section*{Project and Book Identifiers}

\begin{center}
\begin{tabular}{|p{0.16\textwidth}|p{0.76\textwidth}|}
\hline
\textbf{Code} & \textbf{Meaning} \\
\hline

VS
&
Viability System; the abbreviated designation used for the
philosophical system developed in
\textit{The Philosophy of Viability}.
\\
\hline

B01
&
Book 01; the structural identifier of the first book in the
project, \textit{Ontology of Viability}.
\\
\hline

\end{tabular}
\end{center}


% =============================================================================
% 02 — LOGICAL AND PHILOSOPHICAL REGISTER CODES
% =============================================================================

\section*{Logical and Philosophical Register Codes}

The following codes identify categories of entries in the project's
formal registers. They are classification identifiers and should not
be interpreted as establishing the epistemic or ontological status
of an entry.

\begin{center}
\begin{tabular}{|p{0.16\textwidth}|p{0.76\textwidth}|}
\hline
\textbf{Code} & \textbf{Meaning} \\
\hline

AX
&
Axiom; an entry classified in the axioms register.
\\
\hline

DF
&
Definition; an entry classified in the definitions register.
\\
\hline

PR
&
Proposition; an entry classified in the propositions register.
\\
\hline

TH
&
Thesis; an entry classified in the theses register.
\\
\hline

HY
&
Hypothesis; an entry classified in the hypotheses register.
\\
\hline

AR
&
Argument; an entry classified in the arguments register.
\\
\hline

PF
&
Philosophical Foundation; an entry classified in the
philosophical foundations register.
\\
\hline

\end{tabular}
\end{center}


% =============================================================================
% 03 — EDITORIAL AND INTERPRETIVE NOTE
% =============================================================================

\section*{Editorial Note}

The codes listed above serve an organizational and referential
function within the philosophical library.

The distinction between an axiom, a definition, a proposition,
a thesis, a hypothesis, an argument, and a philosophical foundation
is a distinction between register categories. Their inclusion
in this list does not imply that all such categories possess
the same logical function, justificatory status, or philosophical
authority.

The classification of a particular entry must be determined
by its content, function, and formal evaluation within the
relevant register.

The functional or logical classification of an entry must also
be distinguished from its epistemic or justificatory status.
A register code alone does not establish whether an entry is
foundationally adopted, stipulated, argued, conditionally supported,
provisional, hypothetical, unresolved, or superseded.

Likewise, the abbreviation \textbf{VS} identifies the philosophical
system, while \textbf{B01} identifies the first book in the series.
Neither identifier should be treated as a substantive ontological
claim.


% =============================================================================
% END OF ACRONYMS AND STRUCTURAL CODES
% =============================================================================